\section{Blueprint Kuis, UTS, UAS, dan Proyek Akhir}
\label{sec:ch17-blueprint}

\subsection{Komponen dan Bobot}

\begin{table}[H]
\centering
\caption{Blueprint asesmen dan cakupan CPMK}
\label{tab:ch17-blueprint}
\begin{tabularx}{\textwidth}{@{}l l X r@{}}
\toprule
\textbf{Komponen} & \textbf{CPMK} & \textbf{Fokus} & \textbf{Bobot} \\
\midrule
Kuis & 01, 02 & Konsep kompilasi, token, automata & 5\% \\
Tugas Teori & 01--03 & Grammar, automata, parsing & 8\% \\
Praktikum 1 & 02 & Lexer & 8\% \\
Praktikum 2 & 03, 04 & Parser & 8\% \\
Praktikum 3 & 04 & Parser generator dan AST & 7\% \\
Praktikum 4 & 01 & Symbol table dan semantik & 8\% \\
UTS & 01--03 & Teori bahasa dan parsing & 20\% \\
Project Final & 02--04 & Mini compiler front-end & 20\% \\
UAS & 01--04 & Komprehensif & 16\% \\
\bottomrule
\end{tabularx}
\end{table}

\subsection{UTS}

UTS menekankan konsep teori dan analisis:
\begin{enumerate}[label=\arabic*.,leftmargin=*]
  \item klasifikasi tata bahasa dan automata (C2--C4);
  \item derivasi, pohon urai, dan ambiguitas (C3--C4);
  \item FIRST/FOLLOW dan tabel LL(1) (C3--C4);
  \item analisis shift-reduce dan konflik (C4--C5);
  \item soal koding singkat (C3).
\end{enumerate}

\subsection{UAS}

UAS menekankan integrasi dan analisis:
\begin{enumerate}[label=\arabic*.,leftmargin=*]
  \item arsitektur kompilator dan lingkungan run-time (C2--C4);
  \item analisis semantik dan type checking (C3--C4);
  \item validasi AST dan diagnostik (C4--C5);
  \item studi kasus program dengan kesalahan (C4--C5);
  \item refleksi dan desain solusi (C5--C6).
\end{enumerate}

\subsection{Proyek Akhir}

Proyek akhir menilai kemampuan integrasi:
\begin{itemize}[leftmargin=*]
  \item pipeline lexer--parser--AST--type checker--validator;
  \item kualitas kode dan pengujian;
  \item dokumentasi dan demo;
  \item kemampuan menjelaskan keputusan desain.
\end{itemize}

\begin{figure}[H]
\centering
\begin{tikzpicture}[node distance=6mm and 8mm]
  \node[phase, text width=2.6cm] (uts) {UTS\\teori};
  \node[phase, text width=2.6cm, right=of uts] (prak)
        {Praktikum\\aplikasi};
  \node[phase, text width=2.6cm, right=of prak] (proj)
        {Proyek\\integrasi};
  \node[phase, text width=2.6cm, right=of proj] (uas)
        {UAS\\komprehensif};
  \draw[arrow] (uts) -- (prak);
  \draw[arrow] (prak) -- (proj);
  \draw[arrow] (proj) -- (uas);
\end{tikzpicture}
\caption{Alur asesmen dari teori menuju integrasi}
\label{fig:ch17-alur}
\end{figure}

\subsection{Contoh Soal UTS}

\begin{enumerate}[label=\arabic*.,leftmargin=*]
  \item Diberikan tata bahasa, tentukan tipe Chomsky dan jelaskan
        alasannya.
  \item Buat derivasi leftmost dan pohon urai untuk string tertentu.
  \item Hitung FIRST dan FOLLOW, lalu bangun tabel LL(1).
  \item Telusuri shift-reduce dan tandai handle pada setiap langkah.
  \item Jelaskan konflik shift-reduce dan cara menyelesaikannya.
\end{enumerate}

\subsection{Contoh Soal UAS}

\begin{enumerate}[label=\arabic*.,leftmargin=*]
  \item Jelaskan perbedaan one-pass dan multi-pass dengan contoh.
  \item Tentukan tipe ekspresi campuran dan jelaskan promosi.
  \item Diberikan AST, tuliskan aturan validasi yang dilanggar.
  \item Rancang pesan kesalahan untuk tiga kasus berbeda.
  \item Usulkan perbaikan desain untuk program yang gagal validasi.
\end{enumerate}
